\begin{lemma}{Eigenschaften der Erweiterung im normierten Fall}
             {EigenschaftenDerErweiterungImNormiertenFall}
Seien $(V, \norm\cdot_V)$, $(W, \norm\cdot_W)$ normierte $\K$--Vektorräume,
$\varphi \in L(V, W)$, $p \in [1, \infty[$ und $n \in \N$. $V^n$ sei mit
$\norm\cdot_{p, n, V}$ und $W^n$ mit $\norm\cdot_{p, n, W}$ ausgestattet.\\
Dann gilt für alle $v \in V^n$:
    $\norm{\pfeil \varphi n(v)}_{p, n, W}
    \le \norm{\varphi}_{V \rightarrow W}\cdot \norm{v}_{p, n, V}$\\
Insbesondere ist $\pfeil \varphi n \in L(V^n, W^n)$
und $\opnorm{\pfeil \varphi n}{V^n}{W^n} = \opnorm{\varphi} V W$, also
\[\pfeil \cdot n: L(V, W) \rightarrow L(V^n, W^n), 
  \alpha \mapsto \pfeil \alpha n\]
eine lineare Isometrie (bezüglich der obigen Normwahl).
\end{lemma}
\begin{proof}
Nach~\ref{EigenschaftenDerErweiterungImVektorraumfall} ist $\pfeil\varphi n$
linear.\\
Wegen $p \ge 1$ sind $\norm \cdot_{p, n, V}$ und $\norm\cdot_{p,n,W}$ Normen.\\
Sei $v \in V^n$.
Dann gilt\[
       \norm{\pfeil \varphi n (v)}_{p, n, W}^p
   =   \sum_{i=1}^n \norm{\varphi(v_i)}_W^p
   \le \sum_{i=1}^n \opnorm\varphi V W^p \norm{v_i}_V^p
   =   \opnorm \varphi V W ^p \cdot \norm v _{p, n, V}^p\text.\]
Damit ist einerseits die erste Aussage gezeigt und andererseits ist 
$\pfeil \varphi n$ nach bekannten Charakterisierungen der Stetigkeit linearer
Abbildungen bereits stetig mit
$\opnorm {\pfeil \varphi n}{V^n}{W^n} \le \opnorm \varphi V W$.\\
Sei nun $\eps \in \R_{>0}$.\\
Dann existiert ein $v_\eps \in \overline K_V(0,1)$ mit
$\opnorm \varphi V W - \norm{\varphi(v_\eps)}_W< \eps$.\\
Sei
\[v' :\haken n \rightarrow V, i \mapsto \begin{cases}v_\eps & : i=1\\
0 &: \text{ sonst} \end{cases}\text.\]
Dann ist offenbar $\norm {v'}_{p, n, V} = \norm {v_\eps}_V \le 1$ und
$\norm {\pfeil \varphi n (v')}_{p, n, W} = \norm{\varphi(v_\eps)}_W$.\\
Also ist $
 \opnorm{\pfeil \varphi n}{V^n}{W^n} - \norm{\pfeil \varphi n (v')}
 \le \opnorm \varphi V W - \norm{\varphi(v_\eps)} < \eps$.\\
Damit ist $\opnorm \varphi V W
 = \sup\menge{\norm{\pfeil \varphi n (x)}_{p,n,W}}{x \in \overline K_{V^n}(0,1)}
 = \opnorm{\pfeil \varphi n}{V^n}{W^n}$.\\
Wegen $L(V, W) \le \Hom(V, W)$ ist $\pfeil \cdot n$ nach
\ref{EigenschaftenDerErweiterungImVektorraumfall} linear.
Die Isometrieeigenschaft ist genau das \enquote{Insbesondere}.
\end{proof}